\documentclass[aspectratio=169]{beamer}

\usetheme{Madrid}

\usepackage[utf8]{inputenc}
\usepackage[T1]{fontenc}
\usepackage[brazil]{babel}
\usepackage{amsmath,amssymb,amsfonts}
\usepackage{graphicx}
\usepackage{tikz}
\usepackage[most]{tcolorbox}
\usepackage{pgfplots}
\pgfplotsset{compat=1.18}
\usetikzlibrary{calc,arrows.meta}

%-------------------------------------------------
% Cores padrao
%-------------------------------------------------
\definecolor{myblue}{RGB}{0,86,125}
\definecolor{mycyan}{RGB}{0,183,235}
\definecolor{myred}{RGB}{160,30,30}
\definecolor{mygray}{RGB}{245,247,250}
\definecolor{mydark}{RGB}{20,35,45}
\definecolor{mygreen}{RGB}{35,120,80}

\setbeamercolor{structure}{fg=myblue}
\setbeamercolor{frametitle}{bg=myblue,fg=white}
\setbeamercolor{title}{fg=white}
\setbeamertemplate{navigation symbols}{}

%-------------------------------------------------
% Boxes padrao
%-------------------------------------------------
\newtcolorbox{mybox}[1]{
colback=mygray,
colframe=myblue,
title=#1,
fonttitle=\bfseries,
arc=3mm,
boxrule=0.8pt
}

\newtcolorbox{eqbox}{
colback=white,
colframe=mycyan,
arc=3mm,
boxrule=0.8pt
}

\newcommand{\sectioncover}[1]{
\begin{frame}[plain]
\begin{tikzpicture}[remember picture,overlay]
\fill[myblue] (current page.south west) rectangle (current page.north east);
\fill[mycyan,opacity=0.25] (current page.north east) circle (4cm);
\fill[mycyan,opacity=0.18] (current page.south west) circle (5cm);
\end{tikzpicture}
\vfill
\begin{center}
{\color{white}\Huge\bfseries #1}
\end{center}
\vfill
\end{frame}
}

\newcommand{\notas}[1]{%
\begin{tikzpicture}[remember picture,overlay]
\node[anchor=south west,
      xshift=3mm,
      yshift=3mm,
      draw,
      rectangle,
      rounded corners=2pt,
      fill=white,
      font=\tiny,
      inner sep=2pt]
at (current page.south west)
	{notas: #1};
\end{tikzpicture}
}

\newcommand{\beautifulcover}{
{
\begin{frame}[plain]
\begin{tikzpicture}[remember picture,overlay]
\fill[mydark] (current page.south west) rectangle (current page.north east);
\fill[myblue!70] ([xshift=-2.5cm,yshift=1.5cm]current page.south west) circle (4.5cm);
\fill[mycyan!45] ([xshift=2.5cm,yshift=-1.0cm]current page.north east) circle (4.2cm);
\fill[white,opacity=0.05] (current page.center) circle (7.5cm);

\node[align=center,text=white] at (current page.center) {
{\Huge\bfseries Matéria Dielétrica}\\[0.35cm]
{\Large Capítulo 6}\\[0.55cm]
{\large Polarização, campo auxiliar, energia e forças}
};

\node[align=center,text=white!75] at ([yshift=-3.1cm]current page.center) {
Prof. Dr. Rafael Camargo Rodrigues de Lima
};
\end{tikzpicture}
\end{frame}
}
}

\begin{document}

\beautifulcover

%=================================================
\sectioncover{6.1 Introdução}

%-------------------------------------------------
\begin{frame}{O que é um dielétrico?}

\begin{mybox}{Definição operacional}
Um dielétrico é um meio que \textbf{não consegue excluir completamente} um campo
elétrico estático macroscópico de seu interior.
\end{mybox}

\vspace{0.25cm}

Em um condutor perfeito:
\[
\mathbf E_{\rm tot}=0
\quad \text{no interior.}
\]

Em um dielétrico:

\begin{eqbox}
\begin{equation}
\mathbf E_{\rm tot}(\mathbf r)
=
\mathbf E_{\rm ext}(\mathbf r)
+
\mathbf E_{\rm self}(\mathbf r)
\neq 0.
\tag{6.1}
\label{eq:campo-total-dieletrico}
\end{equation}
\end{eqbox}

\begin{alertblock}{Origem física}
Ligações químicas e efeitos quânticos restringem o rearranjo das cargas internas.
\end{alertblock}

\end{frame}

%-------------------------------------------------
\begin{frame}{Condutor versus dielétrico}

\begin{columns}
\column{0.48\textwidth}
\begin{mybox}{Condutor perfeito}
\begin{itemize}
\item Cargas livres se rearranjam sem barreira macroscópica.
\item O campo interno estático é nulo.
\item A carga excedente fica na superfície.
\end{itemize}
\end{mybox}

\[
\mathbf E_{\rm in}=0.
\]

\column{0.48\textwidth}
\begin{mybox}{Dielétrico}
\begin{itemize}
\item Cargas deslocam-se apenas um pouco.
\item O campo interno é reduzido, mas não eliminado.
\item Surgem cargas de polarização.
\end{itemize}
\end{mybox}

\[
\mathbf E_{\rm in}\neq 0.
\]
\end{columns}

\vspace{0.25cm}

\begin{alertblock}{Ideia de continuidade}
Aqui \(\kappa\equiv\epsilon/\epsilon_0\) é a \textbf{constante dielétrica}, ou
permissividade relativa. O limite \(\kappa\to\infty\) de muitos resultados
dielétricos reproduz o comportamento de condutores.
\end{alertblock}

\end{frame}

%=================================================
\sectioncover{6.2 Polarização}

%-------------------------------------------------
\begin{frame}{Carga livre e carga de polarização}

\small

O campo externo é produzido por cargas que não pertencem ao dielétrico. Chamamos essas cargas de
\textbf{cargas livres}.

\begin{mybox}{Densidade total de carga}
A carga que entra nas equações de Maxwell é a soma:
\end{mybox}

\begin{eqbox}
\begin{equation}
\rho(\mathbf r)
=
\rho_f(\mathbf r)
+
\rho_P(\mathbf r).
\tag{6.2}
\label{eq:carga-total-polarizacao}
\end{equation}
\end{eqbox}

\begin{itemize}
\item \(\rho_f\): carga livre, colocada ou controlada externamente.
\item \(\rho_P\): carga de polarização, induzida dentro do material.
\end{itemize}

\begin{alertblock}{Leitura física}
\(\rho_P\) é a carga macroscópica que aparece porque cargas positivas e negativas
se deslocam levemente em sentidos opostos.
\end{alertblock}

\end{frame}

%-------------------------------------------------
\begin{frame}{A função polarização \(\mathbf P(\mathbf r)\)}

\small

Para uma amostra neutra de volume \(V\) e superfície \(S\), a carga de polarização
total deve somar zero:

\begin{eqbox}
\begin{equation}
\int_V d^3r\,\rho_P(\mathbf r)
+
\int_S dS\,\sigma_P(\mathbf r_S)
=0.
\tag{6.3}
\label{eq:neutralidade-polarizacao}
\end{equation}
\end{eqbox}

Zangwill introduz \(\mathbf P(\mathbf r)\) por:

\begin{eqbox}
\begin{align}
\rho_P(\mathbf r)
&=
-\nabla\cdot \mathbf P(\mathbf r),
\qquad \mathbf r\in V,
\tag{6.4}
\\
\sigma_P(\mathbf r_S)
&=
\mathbf P(\mathbf r_S)\cdot\hat{\mathbf n}(\mathbf r_S),
\qquad \mathbf r_S\in S,
\tag{6.5}
\\
\mathbf P(\mathbf r)&=0,
\qquad \mathbf r\notin V.
\tag{6.6}
\end{align}
\end{eqbox}

\notas{pag.1-2}

\end{frame}

%-------------------------------------------------
\begin{frame}{Por que essas definições garantem neutralidade?}

\footnotesize

Substituindo \(\rho_P=-\nabla\cdot\mathbf P\) e
\(\sigma_P=\mathbf P\cdot\hat{\mathbf n}\):

\[
Q_P
=
\int_V d^3r\,\rho_P
+
\int_S dS\,\sigma_P.
\]

\[
Q_P
=
-\int_V d^3r\,\nabla\cdot\mathbf P
+
\int_S dS\,\mathbf P\cdot\hat{\mathbf n}.
\]

Pelo teorema da divergência,

\[
\int_V d^3r\,\nabla\cdot\mathbf P
=
\int_S dS\,\mathbf P\cdot\hat{\mathbf n}.
\]

\begin{eqbox}
\begin{equation}
Q_P=0.
\end{equation}
\end{eqbox}

\begin{alertblock}{Moral}
A polarização transforma uma redistribuição interna de carga em uma combinação
consistente de carga de volume e carga de superfície.
\end{alertblock}

\end{frame}

%-------------------------------------------------
\begin{frame}{Desenho: origem de \(\rho_P\) e \(\sigma_P\)}

\begin{center}
\begin{tikzpicture}[scale=0.95]
\draw[thick,myblue,fill=mygray] (-3,-1.2) rectangle (3,1.2);
\node at (0,-1.55) {\small dielétrico polarizado};

\foreach \x in {-2.4,-1.6,-0.8,0,0.8,1.6,2.4}{
  \draw[->,thick,mygreen] (\x,-0.65) -- (\x+0.35,-0.65);
  \draw[->,thick,mygreen] (\x,0) -- (\x+0.35,0);
  \draw[->,thick,mygreen] (\x,0.65) -- (\x+0.35,0.65);
}

\foreach \y in {-0.8,-0.35,0.1,0.55,1.0}{
  \node at (-3.22,\y) {\(-\)};
  \node[myred] at (3.22,\y) {\(+\)};
}

\node[mygreen] at (0,1.55) {\(\mathbf P\)};
\draw[->,very thick,mygreen] (-0.45,1.45) -- (0.45,1.45);
\node at (-3.8,0) {\(\sigma_P<0\)};
\node at (3.8,0) {\(\sigma_P>0\)};
\end{tikzpicture}
\end{center}

\begin{mybox}{Interpretação}
Se \(\mathbf P\) é uniforme, cargas internas vizinhas se compensam; sobra carga
apenas onde a polarização termina, isto é, na superfície.
\end{mybox}

\end{frame}

%-------------------------------------------------
\begin{frame}{Integral de volume de \(\mathbf P\)}

\small

Queremos mostrar que a integral de \(\mathbf P\) é o momento de dipolo total:

\[
\mathbf p=\int_V d^3r\,\mathbf P(\mathbf r).
\]

Para o componente cartesiano \(k\), use:

\[
\nabla\cdot(r_k\mathbf P)
=
\mathbf P\cdot\nabla r_k
+
r_k\nabla\cdot\mathbf P.
\]

Como \(\nabla r_k=\hat{\mathbf e}_k\),

\[
\mathbf P\cdot\nabla r_k=P_k.
\]

Logo,

\begin{eqbox}
\begin{equation}
P_k
=
\nabla\cdot(r_k\mathbf P)
-
r_k\nabla\cdot\mathbf P.
\tag{6.7}
\label{eq:Pk-identidade}
\end{equation}
\end{eqbox}
\notas{pag. 3-4}
\end{frame}

%-------------------------------------------------
\begin{frame}{Demonstração de \(\int \mathbf P=\mathbf p\)}

\small

Integre (6.7) no volume \(V\):

\[
\int_V d^3r\,P_k
=
\int_V d^3r\,\nabla\cdot(r_k\mathbf P)
-
\int_V d^3r\,r_k\nabla\cdot\mathbf P.
\]

Agora use:

\[
\int_V d^3r\,\nabla\cdot(r_k\mathbf P)
=
\int_S dS\,r_k\,\mathbf P\cdot\hat{\mathbf n}
=
\int_S dS\,r_k\sigma_P,
\]

e

\[
-\nabla\cdot\mathbf P=\rho_P.
\]

\begin{eqbox}
\begin{equation}
\int_V d^3r\,P_k
=
\int_S dS\,r_k\sigma_P
+
\int_V d^3r\,r_k\rho_P.
\tag{6.8}
\label{eq:Pk-dipolo}
\end{equation}
\end{eqbox}

O lado direito é exatamente o componente \(p_k\) do momento de dipolo.

\end{frame}

%-------------------------------------------------
\begin{frame}{Resultado: polarização como densidade de dipolo}

\begin{eqbox}
\begin{equation}
\int_V d^3r\,\mathbf P(\mathbf r)
=
\mathbf p.
\tag{6.9}
\label{eq:P-integral-dipolo}
\end{equation}
\end{eqbox}

\begin{mybox}{Interpretação cuidadosa}
É tentador dizer que \(\mathbf P\) é o ``momento de dipolo por unidade de
volume''. Isso é útil, mas não é uma definição microscópica universal.
\end{mybox}

\begin{alertblock}{Aviso importante}
Em sólidos reais, a polarização não é simplesmente a soma de dipolos moleculares
localizados. Ela depende de correntes de polarização e da estrutura quântica do
estado eletrônico.
\end{alertblock}

\end{frame}

%-------------------------------------------------
\begin{frame}{Modelo de Lorentz}

\small

O modelo de Lorentz escolhe uma pequena célula microscópica \(\Omega\), centrada
no ponto macroscópico \(\mathbf r\), e define:

\begin{eqbox}
\begin{equation}
\mathbf P(\mathbf r)
=
\frac{1}{\Omega}
\int_{\Omega} d^3s\,
\mathbf s\,\rho_{\rm micro}(\mathbf s)
=
\frac{\mathbf p(\mathbf r)}{\Omega}.
\tag{6.10}
\label{eq:lorentz-P}
\end{equation}
\end{eqbox}

\begin{mybox}{Quando funciona bem?}
Gases, líquidos não polares e sólidos moleculares fracamente ligados.
\end{mybox}

\begin{alertblock}{Quando falha?}
Em dielétricos covalentes ou iônicos, a escolha da célula pode mudar o resultado
e correntes entre células podem dominar a polarização real.
\end{alertblock}

\end{frame}

%-------------------------------------------------
\begin{frame}{Teoria moderna da polarização}

\small

Zangwill enfatiza que, em sólidos, \(\mathbf P\) não é determinado apenas pela
densidade de carga estática.

\begin{mybox}{Ideia essencial}
A polarização macroscópica pode ser calculada por uma média temporal de
correntes de polarização:
\end{mybox}

\begin{eqbox}
\begin{equation}
\mathbf P(\mathbf R)
=
\frac{1}{\Omega}
\int_{-\infty}^{t}dt'
\int_{\Omega}d^3s\,
\mathbf j_{P,\rm micro}(\mathbf s,t').
\tag{6.11}
\label{eq:P-corrente}
\end{equation}
\end{eqbox}

\begin{alertblock}{Mensagem}
A polarização mede ``quanto carregamento atravessou'' uma célula ao ligar
adiabaticamente o campo, não apenas onde a carga ficou no final.
\end{alertblock}

\end{frame}

%=================================================
\sectioncover{6.3 Campo produzido por matéria polarizada}

%-------------------------------------------------
\begin{frame}{Potencial da carga de polarização}

\small

Se \(\mathbf P(\mathbf r)\) é dado, então as fontes do campo produzido pela
matéria polarizada são:

\[
\rho_P=-\nabla\cdot\mathbf P,
\qquad
\sigma_P=\mathbf P\cdot\hat{\mathbf n}.
\]

O potencial é:

\begin{eqbox}
\begin{equation}
\varphi_P(\mathbf r)
=
\frac{1}{4\pi\epsilon_0}
\int_V d^3r'\,
\frac{-\nabla'\cdot\mathbf P(\mathbf r')}
{|\mathbf r-\mathbf r'|}
+
\frac{1}{4\pi\epsilon_0}
\int_S dS'\,
\frac{\mathbf P(\mathbf r')\cdot\hat{\mathbf n}'}
{|\mathbf r-\mathbf r'|}.
\tag{6.12}
\label{eq:potencial-polarizacao}
\end{equation}
\end{eqbox}

\end{frame}

%-------------------------------------------------
\begin{frame}{Campo da matéria polarizada}

\small

O campo é \(\mathbf E_P=-\nabla\varphi_P\). Diferenciando em relação ao ponto de
observação \(\mathbf r\):

\begin{eqbox}
\begin{equation}
\mathbf E_P(\mathbf r)
=
\frac{1}{4\pi\epsilon_0}
\int_V d^3r'\,
\big[-\nabla'\cdot\mathbf P(\mathbf r')\big]
\frac{\mathbf r-\mathbf r'}{|\mathbf r-\mathbf r'|^3}
+
\frac{1}{4\pi\epsilon_0}
\int_S dS'\,
\big[\mathbf P(\mathbf r')\cdot\hat{\mathbf n}'\big]
\frac{\mathbf r-\mathbf r'}{|\mathbf r-\mathbf r'|^3}.
\tag{6.13}
\label{eq:campo-polarizacao}
\end{equation}
\end{eqbox}

\begin{alertblock}{Ponto técnico}
Essa forma é bem comportada mesmo para pontos de observação dentro do material.
\end{alertblock}

\end{frame}

%-------------------------------------------------
\begin{frame}{Poisson: amostra uniformemente polarizada}

\small

Se \(\mathbf P\) é uniforme:

\[
\nabla\cdot\mathbf P=0.
\]

Então só resta a integral de superfície:

\[
\mathbf E_P(\mathbf r)
=
\frac{1}{4\pi\epsilon_0}
\int_S dS'\,
(\mathbf P\cdot\hat{\mathbf n}')
\frac{\mathbf r-\mathbf r'}{|\mathbf r-\mathbf r'|^3}.
\]

Use o teorema da divergência em \(\mathbf r'\):

\[
\int_S dS'\,\hat n'_k F(\mathbf r')
=
\int_V d^3r'\,\nabla'_k F(\mathbf r').
\]

Com
\[
F(\mathbf r')=\frac{\mathbf r-\mathbf r'}{|\mathbf r-\mathbf r'|^3}.
\]

\end{frame}

%-------------------------------------------------
\begin{frame}{Fórmula de Poisson}

\small

Como \(\mathbf P\) é constante, podemos tirá-la da integral:

\[
\mathbf E_P(\mathbf r)
=
\mathbf P\cdot\nabla'
\left[
\frac{1}{4\pi\epsilon_0}
\int_V d^3r'\,
\frac{\mathbf r-\mathbf r'}{|\mathbf r-\mathbf r'|^3}
\right].
\]

Mas

\[
\nabla' f(\mathbf r-\mathbf r')=-\nabla f(\mathbf r-\mathbf r').
\]

Defina \(\mathbf E_1(\mathbf r)\) como o campo do mesmo corpo preenchido por uma
densidade uniforme \(\rho=1\). Então:

\begin{eqbox}
\begin{equation}
\mathbf E_P(\mathbf r)
=
-(\mathbf P\cdot\nabla)\mathbf E_1(\mathbf r).
\tag{6.15}
\label{eq:poisson-polarizacao}
\end{equation}
\end{eqbox}

\end{frame}

%-------------------------------------------------
\begin{frame}{Aplicação 6.1: esfera uniformemente polarizada}

\small

Para uma esfera de raio \(R\), primeiro calcule o campo de uma esfera com
densidade uniforme \(\rho=1\):

\begin{eqbox}
\begin{equation}
\mathbf E_1(\mathbf r)
=
\begin{cases}
\dfrac{\mathbf r}{3\epsilon_0}, & r<R,\\[0.25cm]
\dfrac{V}{4\pi\epsilon_0}\dfrac{\mathbf r}{r^3}, & r>R.
\end{cases}
\tag{6.16}
\label{eq:campo-esfera-rho1}
\end{equation}
\end{eqbox}

Agora aplique
\[
\mathbf E_P=-(\mathbf P\cdot\nabla)\mathbf E_1.
\]

\begin{mybox}{Interior}
Como \(\nabla_i r_j=\delta_{ij}\),
\[
\mathbf E_P(r<R)=-\frac{\mathbf P}{3\epsilon_0}.
\]
\end{mybox}

\end{frame}

%-------------------------------------------------
\begin{frame}{Campo da esfera polarizada}

\small

O resultado completo é:

\begin{eqbox}
\begin{equation}
\mathbf E_P(\mathbf r)
=
\begin{cases}
-\dfrac{\mathbf P}{3\epsilon_0},
& r<R,\\[0.35cm]
\dfrac{V}{4\pi\epsilon_0}
\dfrac{3(\hat{\mathbf r}\cdot\mathbf P)\hat{\mathbf r}-\mathbf P}{r^3},
& r>R.
\end{cases}
\tag{6.17}
\label{eq:campo-esfera-polarizada}
\end{equation}
\end{eqbox}

\begin{alertblock}{Leitura}
Fora da esfera, o campo é exatamente o de um dipolo pontual
\[
\mathbf p=V\mathbf P.
\]
Dentro, o campo é uniforme e antiparalelo a \(\mathbf P\).
\end{alertblock}

\end{frame}

%-------------------------------------------------
\begin{frame}{Método alternativo: carga superficial}

\small

Se \(\mathbf P=P\hat{\mathbf z}\), então:

\[
\sigma_P=\mathbf P\cdot\hat{\mathbf r}=P\cos\theta.
\tag{6.18}
\]

Esse é o mesmo tipo de distribuição superficial que apareceu no capítulo de
multipolos:

\[
\sigma(\theta)=\sigma_0\cos\theta.
\]

O potencial é:

\begin{eqbox}
\begin{equation}
\varphi_P(r,\theta)
=
\frac{P}{3\epsilon_0}
\begin{cases}
z, & r<R,\\[0.2cm]
\dfrac{R^3}{r^2}\cos\theta, & r>R.
\end{cases}
\tag{6.19}
\label{eq:potencial-esfera-polarizada}
\end{equation}
\end{eqbox}

\end{frame}

%-------------------------------------------------
\begin{frame}{Matéria polarizada como coleção de dipolos}

\small

Queremos reescrever (6.12) sem uma integral de superfície explícita. Pelo teorema
da divergência:

\[
\int_S dS'\,\frac{\mathbf P\cdot\hat{\mathbf n}'}{|\mathbf r-\mathbf r'|}
=
\int_V d^3r'\,
\nabla'\cdot
\left[
\frac{\mathbf P(\mathbf r')}{|\mathbf r-\mathbf r'|}
\right].
\]

Expanda a divergência:

\[
\nabla'\cdot
\left(
\frac{\mathbf P}{|\mathbf r-\mathbf r'|}
\right)
=
\frac{\nabla'\cdot\mathbf P}{|\mathbf r-\mathbf r'|}
+
\mathbf P\cdot\nabla'
\frac{1}{|\mathbf r-\mathbf r'|}.
\]

O primeiro termo cancela o termo de volume original em (6.12).

\end{frame}

%-------------------------------------------------
\begin{frame}{Resultado fundamental}

\small

Após a cancelamento, fica:

\begin{eqbox}
\begin{equation}
\varphi_P(\mathbf r)
=
\frac{1}{4\pi\epsilon_0}
\int_V d^3r'\,
\mathbf P(\mathbf r')\cdot
\nabla'
\frac{1}{|\mathbf r-\mathbf r'|}.
\tag{6.22}
\label{eq:potencial-dipolos-continuos}
\end{equation}
\end{eqbox}

\begin{mybox}{Comparação com um dipolo pontual}
O potencial de um dipolo \(\mathbf p\) em \(\mathbf r'\) tem a forma
\[
\varphi_{\rm dip}
=
\frac{1}{4\pi\epsilon_0}
\mathbf p\cdot\nabla'
\frac{1}{|\mathbf r-\mathbf r'|}.
\]
\end{mybox}

\begin{alertblock}{Conclusão}
Matéria polarizada é eletrostaticamente equivalente a uma distribuição contínua
de dipolos pontuais \(d\mathbf p=\mathbf P\,d^3r\).
\end{alertblock}

\end{frame}

%=================================================
\sectioncover{6.4 O campo total e o campo auxiliar \(\mathbf D\)}

%-------------------------------------------------
\begin{frame}{Lei de Gauss em matéria}

\small

A carga total é a soma de carga livre e carga de polarização:

\[
\rho=\rho_f+\rho_P.
\]

Como
\[
\rho_P=-\nabla\cdot\mathbf P,
\]
temos:

\begin{eqbox}
\begin{equation}
\epsilon_0\nabla\cdot\mathbf E
=
\rho_f-\nabla\cdot\mathbf P.
\tag{6.24}
\label{eq:gauss-materia}
\end{equation}
\end{eqbox}

Reorganizando:

\[
\nabla\cdot(\epsilon_0\mathbf E+\mathbf P)=\rho_f.
\]

\end{frame}

%-------------------------------------------------
\begin{frame}{Definição do campo auxiliar \(\mathbf D\)}

\begin{eqbox}
\begin{equation}
\mathbf D
=
\epsilon_0\mathbf E+\mathbf P.
\tag{6.25}
\label{eq:D-def}
\end{equation}
\end{eqbox}

\begin{mybox}{Equações eletrostáticas em matéria}
\[
\nabla\times\mathbf E=0,
\tag{6.26}
\]
\[
\nabla\cdot\mathbf D=\rho_f,
\tag{6.27}
\]
\[
\nabla\times\mathbf D=\nabla\times\mathbf P.
\tag{6.28}
\]
\end{mybox}

\begin{alertblock}{Cuidado}
\(\mathbf D\) não é produzido apenas por carga livre. Variações espaciais de
\(\mathbf P\) também entram em sua estrutura.
\end{alertblock}

\end{frame}

%-------------------------------------------------
\begin{frame}{Forma integral de Gauss para \(\mathbf D\)}

\begin{eqbox}
\begin{equation}
\oint_S d\mathbf S\cdot\mathbf D
=
Q_f.
\tag{6.29}
\label{eq:gauss-D-integral}
\end{equation}
\end{eqbox}

\begin{mybox}{Utilidade}
Quando a simetria é boa, essa equação calcula \(\mathbf D\) diretamente a partir
da carga livre.
\end{mybox}

\begin{alertblock}{Mas não confundir}
Ela não diz que a carga livre é a única causa física de \(\mathbf D\). O campo
auxiliar é uma combinação de \(\epsilon_0\mathbf E\) e \(\mathbf P\).
\end{alertblock}

\end{frame}

%-------------------------------------------------
\begin{frame}{Condições de contorno para \(\mathbf D\)}

\small

Em uma interface entre dois meios, com densidade superficial de carga livre
\(\sigma_f\):

\begin{eqbox}
\begin{align}
\hat{\mathbf n}_2\cdot(\mathbf D_1-\mathbf D_2)
&=
\sigma_f,
\tag{6.31}
\label{eq:bc-D-normal}
\\
\hat{\mathbf n}_2\times(\mathbf D_1-\mathbf D_2)
&=
\hat{\mathbf n}_2\times(\mathbf P_1-\mathbf P_2).
\tag{6.32}
\label{eq:bc-D-tangencial}
\end{align}
\end{eqbox}

Como \(\mathbf D=\epsilon_0\mathbf E+\mathbf P\), a segunda condição equivale à
continuidade da componente tangencial de \(\mathbf E\):

\begin{eqbox}
\begin{equation}
\hat{\mathbf n}_2\times(\mathbf E_1-\mathbf E_2)=0.
\tag{6.33}
\label{eq:bc-E-tangencial}
\end{equation}
\end{eqbox}

\end{frame}

%-------------------------------------------------
\begin{frame}{Relações constitutivas}

\small

As equações macroscópicas ainda não fecham o problema. Precisamos relacionar
\(\mathbf P\) com \(\mathbf E\).

\begin{mybox}{Expansão fenomenológica}
\[
P_i
=
\epsilon_0\chi_{ij}E_j
+
\epsilon_0\chi^{(2)}_{ijk}E_jE_k+\cdots.
\tag{6.34}
\]
\end{mybox}

\begin{itemize}
\item Termo linear: resposta usual em campos fracos.
\item Tensores: permitem anisotropia.
\item Termos de ordem superior: óptica não linear e campos intensos.
\end{itemize}

\begin{alertblock}{No restante}
O caso mais importante será o dielétrico simples: linear, isotrópico e homogêneo
por regiões.
\end{alertblock}

\end{frame}

%=================================================
\sectioncover{6.5 Matéria dielétrica simples}

%-------------------------------------------------
\begin{frame}{Dielétrico simples}

\begin{mybox}{Constitutiva linear e isotrópica}
\[
\mathbf P=\epsilon_0\chi\,\mathbf E.
\tag{6.35}
\]
\end{mybox}

Definimos a permissividade \(\epsilon\) e a constante dielétrica \(\kappa\):

\begin{eqbox}
\begin{equation}
\mathbf P
=
(\epsilon-\epsilon_0)\mathbf E
=
\epsilon_0(\kappa-1)\mathbf E.
\tag{6.36}
\label{eq:P-simple}
\end{equation}
\end{eqbox}

Então:

\begin{eqbox}
\begin{equation}
\mathbf D
=
\epsilon\mathbf E
=
\kappa\epsilon_0\mathbf E.
\tag{6.37}
\label{eq:D-simple}
\end{equation}
\end{eqbox}

\end{frame}

%-------------------------------------------------
\begin{frame}{Equações em uma região de \(\epsilon\) constante}

\footnotesize

Use \(\mathbf D=\epsilon\mathbf E\) em \(\nabla\cdot\mathbf D=\rho_f\):

\[
\nabla\cdot(\epsilon\mathbf E)=\rho_f.
\]

Se \(\epsilon\) varia no espaço:

\begin{eqbox}
\begin{equation}
\epsilon\nabla\cdot\mathbf E
+
\mathbf E\cdot\nabla\epsilon
=
\rho_f.
\tag{6.38}
\end{equation}
\end{eqbox}

Se \(\epsilon\) é constante em cada região:

\begin{eqbox}
\begin{equation}
\nabla\cdot\mathbf E
=
\frac{\rho_f}{\epsilon}.
\tag{6.39}
\end{equation}
\end{eqbox}

Como \(\mathbf E=-\nabla\varphi\):

\begin{eqbox}
\begin{equation}
\nabla^2\varphi=-\frac{\rho_f}{\epsilon}.
\tag{6.40}
\end{equation}
\end{eqbox}

\end{frame}

%-------------------------------------------------
\begin{frame}{Carga de polarização em dielétrico simples}

\small

De \(\mathbf P=\epsilon_0(\kappa-1)\mathbf E\):

\[
\rho_P=-\nabla\cdot\mathbf P
=
-\epsilon_0(\kappa-1)\nabla\cdot\mathbf E.
\]

Em região homogênea, \(\nabla\cdot\mathbf E=\rho_f/\epsilon=\rho_f/(\kappa\epsilon_0)\).

\begin{eqbox}
\begin{equation}
\rho_P
=
\left(\frac{1}{\kappa}-1\right)\rho_f.
\tag{6.43}
\label{eq:rhoP-simple}
\end{equation}
\end{eqbox}

Na superfície:

\begin{eqbox}
\begin{equation}
\sigma_P
=
\mathbf P\cdot\hat{\mathbf n}
=
\epsilon_0(\kappa-1)\mathbf E\cdot\hat{\mathbf n}.
\tag{6.44}
\label{eq:sigmaP-simple}
\end{equation}
\end{eqbox}

\end{frame}

%-------------------------------------------------
\begin{frame}{Resposta a uma carga pontual livre}

\small

Para uma carga livre \(q_0\) em um meio simples infinito:

\[
\rho_f(\mathbf r)=q_0\delta(\mathbf r).
\]

Pela equação (6.43):

\begin{eqbox}
\begin{equation}
\rho_P(\mathbf r)
=
-\left(1-\frac{1}{\kappa}\right)
q_0\delta(\mathbf r).
\tag{6.45}
\label{eq:screening-carga-ponto}
\end{equation}
\end{eqbox}

Logo, a carga efetiva macroscópica é:

\[
q_{\rm eff}
=
q_0+q_P
=
\frac{q_0}{\kappa}.
\]

\begin{alertblock}{Screening}
O dielétrico reduz o campo de Coulomb por um fator \(\kappa\).
\end{alertblock}

\end{frame}

%-------------------------------------------------
\begin{frame}{Campo e potencial de uma carga em meio simples}

\small

Pela lei de Gauss para \(\mathbf D\):

\[
\oint d\mathbf S\cdot\mathbf D=q_0.
\]

Assim,

\[
\mathbf D(r)=\frac{q_0}{4\pi r^2}\hat{\mathbf r}.
\]

Como \(\mathbf D=\epsilon\mathbf E=\kappa\epsilon_0\mathbf E\):

\begin{eqbox}
\begin{equation}
\mathbf E(r)
=
\frac{1}{4\pi\kappa\epsilon_0}
\frac{q_0}{r^2}\hat{\mathbf r}.
\tag{6.46}
\end{equation}
\end{eqbox}

\begin{eqbox}
\begin{equation}
\varphi(r)
=
\frac{1}{4\pi\kappa\epsilon_0}
\frac{q_0}{r}.
\tag{6.47}
\end{equation}
\end{eqbox}

\end{frame}

%-------------------------------------------------
\begin{frame}{Capacitor: carga fixa}

\footnotesize

Considere placas paralelas com carga livre fixa \(\pm Q_f\), área \(A\) e
separação \(d\). Então:

\[
D=\sigma_f=\frac{Q_f}{A}.
\]

Com dielétrico simples entre as placas:

\[
E=\frac{D}{\epsilon}=\frac{\sigma_f}{\kappa\epsilon_0}.
\]

Logo:

\begin{eqbox}
\begin{equation}
\Delta\varphi
=
Ed
=
\frac{\sigma_f d}{\kappa\epsilon_0}.
\tag{6.48}
\end{equation}
\end{eqbox}

Como \(C=Q_f/\Delta\varphi\):

\begin{eqbox}
\begin{equation}
C=\kappa C_0,
\qquad
C_0=\frac{\epsilon_0 A}{d}.
\tag{6.52}
\end{equation}
\end{eqbox}

\end{frame}

%-------------------------------------------------
\begin{frame}{Interface simples: carga no centro de uma esfera}

\small

Uma carga \(q\) está no centro de uma esfera de raio \(R\) e permissividade
\(\epsilon_1=\kappa_1\epsilon_0\). Fora dela há meio \(\epsilon_2=\kappa_2\epsilon_0\).

\begin{center}
\begin{tikzpicture}[scale=0.9]
\draw[thick,myblue,fill=mygray] (0,0) circle (1.3);
\node[myred] at (0,0) {\Large \(q\)};
\draw[->,thick] (0,0) -- (1.3,0) node[midway,above] {\(R\)};
\node at (0,-1.75) {\(\kappa_1\)};
\node at (2.15,1.2) {\(\kappa_2\)};
\foreach \a in {0,45,...,315}{
  \draw[->,thin] ({0.25*cos(\a)},{0.25*sin(\a)}) -- ({2.0*cos(\a)},{2.0*sin(\a)});
}
\end{tikzpicture}
\end{center}

Pela simetria esférica:

\begin{eqbox}
\begin{equation}
\mathbf D(r)
=
\frac{q}{4\pi r^2}\hat{\mathbf r},
\qquad r>0.
\tag{6.53}
\end{equation}
\end{eqbox}

\end{frame}

%-------------------------------------------------
\begin{frame}{Carga de polarização na interface esférica}

\small

Em cada lado da interface:

\[
\mathbf E_i=\frac{\mathbf D}{\epsilon_i},
\qquad
\mathbf P_i=(\epsilon_i-\epsilon_0)\mathbf E_i.
\]

Logo,

\[
\mathbf P_i
=
\left(1-\frac{1}{\kappa_i}\right)\mathbf D.
\]

A carga superficial de polarização recebe contribuição dos dois lados:

\[
\sigma_P
=
\hat{\mathbf r}\cdot(\mathbf P_1-\mathbf P_2).
\]

Assim:

\begin{eqbox}
\begin{equation}
\sigma_P(R)
=
\frac{q}{4\pi R^2}
\left(
\frac{1}{\kappa_2}
-
\frac{1}{\kappa_1}
\right).
\tag{6.58}
\label{eq:sigma-interface-esferica}
\end{equation}
\end{eqbox}

\end{frame}

%-------------------------------------------------
\begin{frame}{Exemplo 6.2: carga perto de uma interface plana}

\small

Uma carga \(q\) está em \(z=-d\). O meio à esquerda tem \(\kappa_L\) e o meio à
direita tem \(\kappa_R\). A interface é o plano \(z=0\).

\begin{center}
\begin{tikzpicture}[scale=0.82]
\fill[mygray] (0,-1.4) rectangle (3,1.4);
\draw[thick,myblue] (0,-1.4) -- (0,1.4);
\node at (-1.5,1.15) {\(\kappa_L\)};
\node at (1.5,1.15) {\(\kappa_R\)};
\node[myred] at (-1.35,0) {\Large \(q\)};
\draw[dashed] (-1.35,0) -- (0,0);
\node at (-0.65,-0.25) {\(d\)};
\node at (0.55,-0.25) {\(\mathbf r_0=(x,y,0)\)};
\end{tikzpicture}
\end{center}

\begin{mybox}{Resultado}
A carga de polarização induzida na interface é
\[
\sigma_P(x,y)
=
\frac{1}{2\pi}
\frac{\kappa_L-\kappa_R}{\kappa_L(\kappa_L+\kappa_R)}
\frac{qd}{(x^2+y^2+d^2)^{3/2}}.
\]
\end{mybox}

\end{frame}

%-------------------------------------------------
\begin{frame}{Exemplo 6.2: de onde vem a fórmula?}

\scriptsize

No ponto \(\mathbf r_0=(x,y,0)\), escrevemos:

\[
\sigma_P=\hat{\mathbf z}\cdot(\mathbf P_L-\mathbf P_R)
=
\epsilon_0(\chi_L\mathbf E_L-\chi_R\mathbf E_R)\cdot\hat{\mathbf z}.
\]

Perto da superfície, o campo da própria folha de carga tem salto:

\[
\epsilon_0 E_{Lz}
=
\frac{1}{4\pi\kappa_L}\frac{qd}{r^3}
-\frac{\sigma_P}{2},
\qquad
\epsilon_0 E_{Rz}
=
\frac{1}{4\pi\kappa_L}\frac{qd}{r^3}
+\frac{\sigma_P}{2},
\]

onde \(r^2=x^2+y^2+d^2\). Como não há carga livre na interface:

\[
\kappa_L E_{Lz}=\kappa_R E_{Rz}.
\]

Resolvendo a equação linear resultante para \(\sigma_P\), obtemos:

\begin{eqbox}
\[
\sigma_P
=
\frac{1}{2\pi}
\frac{\kappa_L-\kappa_R}{\kappa_L(\kappa_L+\kappa_R)}
\frac{qd}{r^3}.
\]
\end{eqbox}

\end{frame}

%-------------------------------------------------
\begin{frame}{Condições de contorno em potencial}

\small

Como \(\mathbf E=-\nabla\varphi\), a continuidade da componente tangencial de
\(\mathbf E\) implica:

\begin{eqbox}
\begin{equation}
\varphi_1=\varphi_2
\qquad \text{na interface.}
\tag{6.70}
\label{eq:bc-potencial-continuo}
\end{equation}
\end{eqbox}

Se há carga livre superficial \(\sigma_f\):

\begin{eqbox}
\begin{equation}
\epsilon_1\frac{\partial\varphi_1}{\partial n}
-
\epsilon_2\frac{\partial\varphi_2}{\partial n}
=
-\sigma_f.
\tag{6.71}
\label{eq:bc-potencial-normal}
\end{equation}
\end{eqbox}

Sem carga livre na interface:

\begin{eqbox}
\begin{equation}
\epsilon_1\frac{\partial\varphi_1}{\partial n}
=
\epsilon_2\frac{\partial\varphi_2}{\partial n}.
\tag{6.73}
\end{equation}
\end{eqbox}

\end{frame}

%-------------------------------------------------
\begin{frame}{Capacitor: potencial fixo}

\small

Agora a bateria mantém \(V\) fixo entre as placas. Portanto:

\begin{eqbox}
\begin{equation}
V=\int_0^d d\ell\,E
\quad\Rightarrow\quad
E=\frac{V}{d}.
\tag{6.62--6.63}
\end{equation}
\end{eqbox}

\begin{mybox}{Diferença em relação à carga fixa}
O campo \(E\) não é reduzido quando o dielétrico é inserido, porque a bateria
fornece carga livre extra às placas.
\end{mybox}

\end{frame}

%-------------------------------------------------
\begin{frame}{Capacitor a potencial fixo: carga livre}

\small

A carga livre na placa passa a ser:

\[
\sigma_f=D=\epsilon E=\kappa\epsilon_0\frac{V}{d}.
\tag{6.65}
\]

Logo, novamente:

\begin{eqbox}
\[
C=\frac{\sigma_f A}{V}
=
\kappa C_0.
\tag{6.66}
\]
\end{eqbox}

\end{frame}

%-------------------------------------------------
\begin{frame}{Teoria de potencial para dielétricos simples}

\small

Em uma região de permissividade constante:

\begin{eqbox}
\begin{equation}
\epsilon_0\nabla^2\varphi
=
-\frac{\rho_f}{\kappa}.
\tag{6.69}
\end{equation}
\end{eqbox}

Se não há carga livre:

\begin{eqbox}
\begin{equation}
\nabla^2\varphi=0.
\tag{6.72}
\end{equation}
\end{eqbox}

\begin{mybox}{Receita}
Resolver problemas dielétricos simples é resolver Laplace/Poisson em cada
região e aplicar:
\[
\varphi_1=\varphi_2,
\qquad
\epsilon_1\partial_n\varphi_1=\epsilon_2\partial_n\varphi_2
\]
quando não há carga livre na interface.
\end{mybox}

\end{frame}

%-------------------------------------------------
\begin{frame}{Exemplo 6.3: cavidade esférica em campo uniforme}

\small

Uma cavidade esférica de raio \(R\) é escavada em um meio infinito de constante
\(\kappa\). Um campo uniforme \(\mathbf E_0=E_0\hat{\mathbf z}\) é mantido por
placas distantes.

\begin{mybox}{Forma dos potenciais}
Pela simetria axial e por \(\nabla^2\varphi=0\):
\[
\varphi_{\rm in}=Ar\cos\theta,
\]
\[
\varphi_{\rm out}
=
\left(-E_0r+\frac{B}{r^2}\right)\cos\theta.
\]
\end{mybox}

O termo \(-E_0r\cos\theta\) é o potencial do campo aplicado distante. O termo
\(B\cos\theta/r^2\) é a resposta dipolar induzida pela cavidade.

\end{frame}

%-------------------------------------------------
\begin{frame}{Exemplo 6.3: determinando \(A\) e \(B\)}

\small

Em \(r=R\), impomos continuidade do potencial:

\[
AR=-E_0R+\frac{B}{R^2}.
\tag{1}
\]

Sem carga livre na interface:

\[
\epsilon_0
\left.
\frac{\partial\varphi_{\rm in}}{\partial r}
\right|_{R}
=
\kappa\epsilon_0
\left.
\frac{\partial\varphi_{\rm out}}{\partial r}
\right|_{R}.
\tag{2}
\]

Como
\[
\partial_r\varphi_{\rm in}=A\cos\theta,
\]
\[
\partial_r\varphi_{\rm out}
=
\left(-E_0-\frac{2B}{r^3}\right)\cos\theta,
\]
o sistema linear dá:

\begin{eqbox}
\[
A=-\frac{3\kappa}{2\kappa+1}E_0,
\qquad
B=-\frac{\kappa-1}{2\kappa+1}E_0R^3.
\]
\end{eqbox}

\end{frame}

%-------------------------------------------------
\begin{frame}{Exemplo 6.3: campo dentro da cavidade}

\small

No interior da cavidade:

\[
\varphi_{\rm in}=Ar\cos\theta=Az.
\]

Logo:

\[
\mathbf E_{\rm in}
=
-\nabla\varphi_{\rm in}
=
-A\hat{\mathbf z}.
\]

Usando o valor de \(A\):

\begin{eqbox}
\begin{equation}
\mathbf E_{\rm in}
=
\frac{3\kappa}{2\kappa+1}
E_0\hat{\mathbf z}.
\end{equation}
\end{eqbox}

\begin{alertblock}{Interpretação}
O campo dentro da cavidade é maior que \(E_0\) para \(\kappa>1\). A carga de
polarização na parede da cavidade reforça o campo no interior vazio.
\end{alertblock}

\end{frame}

%-------------------------------------------------
\begin{frame}{Refração de linhas de campo}

\small

Em uma interface plana sem carga livre:

\[
E_{1\parallel}=E_{2\parallel},
\qquad
\epsilon_1 E_{1\perp}=\epsilon_2 E_{2\perp}.
\]

Se \(\theta_i\) é o ângulo entre \(\mathbf E_i\) e a normal:

\[
\tan\theta_i=\frac{E_{i\parallel}}{E_{i\perp}}.
\]

Dividindo as duas relações:

\begin{eqbox}
\begin{equation}
\frac{\tan\theta_1}{\tan\theta_2}
=
\frac{\epsilon_1}{\epsilon_2}
=
\frac{\kappa_1}{\kappa_2}.
\tag{6.68}
\label{eq:refracao-linhas-E}
\end{equation}
\end{eqbox}

\begin{alertblock}{Leitura}
Linhas de \(\mathbf E\) dobram ao cruzar uma interface dielétrica.
\end{alertblock}

\end{frame}

%=================================================
\sectioncover{6.6 A física da constante dielétrica}

%-------------------------------------------------
\begin{frame}{Polarizabilidade elétrica}

\small

Para um objeto pequeno em um campo aproximadamente uniforme, definimos:

\begin{eqbox}
\begin{equation}
\mathbf p=\alpha\epsilon_0\mathbf E.
\tag{6.75}
\label{eq:polarizabilidade-micro}
\end{equation}
\end{eqbox}

\begin{mybox}{Interpretação}
\(\alpha\) mede o quão facilmente a nuvem eletrônica e os núcleos se deformam
sob ação do campo.
\end{mybox}

Para uma esfera condutora, do capítulo 5:

\[
\alpha=4\pi R^3.
\]

Para átomos e moléculas, \(\alpha\) é uma propriedade microscópica medida ou
calculada.

\end{frame}

%-------------------------------------------------
\begin{frame}{Estimativa ingênua da constante dielétrica}

\footnotesize

Se há \(n\) moléculas por unidade de volume e cada uma adquire

\[
\mathbf p=\alpha\epsilon_0\mathbf E,
\]

então:

\[
\mathbf P=n\mathbf p
=
n\alpha\epsilon_0\mathbf E.
\]

Comparando com

\[
\mathbf P=\epsilon_0(\kappa-1)\mathbf E,
\]

obtemos:

\begin{eqbox}
\begin{equation}
\kappa\simeq 1+n\alpha.
\tag{6.79}
\label{eq:kappa-ingenua}
\end{equation}
\end{eqbox}

\begin{alertblock}{Limitação}
Essa fórmula ignora que cada molécula responde ao campo local, não exatamente ao
campo macroscópico médio.
\end{alertblock}

\end{frame}

%-------------------------------------------------
\begin{frame}{Campo local de Lorentz}

\small

O campo que polariza uma molécula é melhor aproximado por:

\[
\mathbf E_{\rm local}
=
\mathbf E
-
\mathbf E_{\rm self}.
\]

Para uma cavidade esférica dentro de um meio uniformemente polarizado:

\[
\mathbf E_{\rm self}
=
-\frac{\mathbf P}{3\epsilon_0}.
\]

Logo:

\begin{eqbox}
\begin{equation}
\mathbf E_{\rm local}
=
\mathbf E
+
\frac{\mathbf P}{3\epsilon_0}.
\tag{6.80--6.81}
\label{eq:campo-local-lorentz}
\end{equation}
\end{eqbox}

\begin{mybox}{Comentário}
O termo \(\mathbf P/3\epsilon_0\) é a correção de Lorentz para o campo de cavidade.
\end{mybox}

\end{frame}

%-------------------------------------------------
\begin{frame}{Clausius-Mossotti: dedução}

\small

Cada molécula tem:

\[
\mathbf p=\alpha\epsilon_0\mathbf E_{\rm local}.
\]

Como \(\mathbf P=n\mathbf p\):

\[
\mathbf P
=
n\alpha\epsilon_0
\left(
\mathbf E+\frac{\mathbf P}{3\epsilon_0}
\right).
\]

Agrupando os termos em \(\mathbf P\):

\[
\left(1-\frac{n\alpha}{3}\right)\mathbf P
=
n\alpha\epsilon_0\mathbf E.
\]

\begin{eqbox}
\begin{equation}
\mathbf P
=
\frac{n\alpha\epsilon_0}{1-n\alpha/3}\mathbf E.
\tag{6.82}
\end{equation}
\end{eqbox}

\end{frame}

%-------------------------------------------------
\begin{frame}{Fórmula de Clausius-Mossotti}

\footnotesize

Compare

\[
\mathbf P
=
\frac{n\alpha\epsilon_0}{1-n\alpha/3}\mathbf E
\]

com

\[
\mathbf P=\epsilon_0(\kappa-1)\mathbf E.
\]

Então:

\[
\kappa-1
=
\frac{n\alpha}{1-n\alpha/3}.
\tag{6.83}
\]

Reorganizando:

\begin{eqbox}
\begin{equation}
\frac{\kappa-1}{\kappa+2}
=
\frac{n\alpha}{3}.
\tag{6.84}
\label{eq:clausius-mossotti}
\end{equation}
\end{eqbox}

\begin{alertblock}{Validade}
Boa para gases não polares e alguns líquidos simples; ruim para muitos sólidos
e líquidos polares.
\end{alertblock}

\end{frame}

%-------------------------------------------------
\begin{frame}{O limite \(\kappa\to\infty\)}

\small

\begin{mybox}{Ideia}
Quando \(\kappa\) cresce muito, o campo macroscópico no interior de um dielétrico
simples tende a ser fortemente reduzido.
\end{mybox}

Em muitos problemas sem carga livre em contato direto com o dielétrico:

\[
\kappa\to\infty
\quad \Longrightarrow \quad
\mathbf E_{\rm in}\to 0.
\]

\begin{alertblock}{Conexão}
Esse é o comportamento de um condutor perfeito.
\end{alertblock}

\begin{mybox}{Cuidado}
Se há carga livre embutida no dielétrico, o limite pode ser singular e não
equivaler simplesmente a um condutor.
\end{mybox}

\end{frame}

%=================================================
\sectioncover{6.7 Energia de matéria dielétrica}

%-------------------------------------------------
\begin{frame}{Trabalho para estabelecer o campo}

\small

Considere um condutor carregado dentro de um dielétrico. Ao trazer uma carga
\(\delta Q\) do infinito até o condutor, fazemos trabalho:

\[
\delta U_E=\varphi\,\delta Q.
\]

Pela lei de Gauss:

\[
\delta Q=\int_S d\mathbf S\cdot\delta\mathbf D.
\]

Então:

\[
\delta U_E
=
\int_S d\mathbf S\cdot\delta\mathbf D\,\varphi.
\]

Aplicando o teorema da divergência no volume externo ao condutor:

\[
\delta U_E
=
-\int_V d^3r\,\nabla\cdot(\varphi\,\delta\mathbf D).
\]

\end{frame}

%-------------------------------------------------
\begin{frame}{Variação da energia}

\small

Expanda:

\[
\nabla\cdot(\varphi\,\delta\mathbf D)
=
\nabla\varphi\cdot\delta\mathbf D
+
\varphi\,\nabla\cdot\delta\mathbf D.
\]

No dielétrico sem carga livre no volume:

\[
\nabla\cdot\delta\mathbf D=0.
\]

Como \(\mathbf E=-\nabla\varphi\):

\begin{eqbox}
\begin{equation}
\delta U_E
=
\int d^3r\,\mathbf E\cdot\delta\mathbf D.
\tag{6.87}
\label{eq:delta-U-dieletrico}
\end{equation}
\end{eqbox}

\begin{alertblock}{Resultado chave}
A energia é uma função natural de \(\mathbf D\), isto é, da carga livre que
determina \(\nabla\cdot\mathbf D=\rho_f\).
\end{alertblock}

\end{frame}

%-------------------------------------------------
\begin{frame}{Energia total geral}

\small

Para construir o campo desde zero até o valor final \(\mathbf D\), integramos:

\begin{eqbox}
\begin{equation}
U_E
=
\int d^3r
\int_0^{\mathbf D}
\mathbf E\cdot\delta\mathbf D.
\tag{6.91}
\label{eq:energia-geral-dieletrico}
\end{equation}
\end{eqbox}

\begin{mybox}{Observação importante}
Não podemos calcular essa integral sem uma relação constitutiva que diga como
\(\mathbf E\) depende de \(\mathbf D\).
\end{mybox}

\begin{alertblock}{Analogia}
Em condutores isolados, \(U_E\) é natural nas cargas \(Q_k\). Em dielétricos,
\(U_E\) é natural em \(\mathbf D\), que é fixado pela carga livre.
\end{alertblock}

\end{frame}

%-------------------------------------------------
\begin{frame}{Energia em dielétrico simples}

\small

Para um dielétrico simples:

\[
\mathbf D=\epsilon\mathbf E,
\qquad
\mathbf E=\frac{\mathbf D}{\epsilon}.
\]

Então:

\[
\int_0^{\mathbf D}\mathbf E\cdot\delta\mathbf D
=
\int_0^{\mathbf D}
\frac{\mathbf D'}{\epsilon}\cdot\delta\mathbf D'
=
\frac{|\mathbf D|^2}{2\epsilon}.
\]

\begin{eqbox}
\begin{equation}
U_E[\mathbf D]
=
\frac{1}{2}
\int d^3r\,\frac{|\mathbf D|^2}{\epsilon}
=
\frac{1}{2}\int d^3r\,\mathbf E\cdot\mathbf D.
\tag{6.94}
\label{eq:energia-simple}
\end{equation}
\end{eqbox}

\end{frame}

%-------------------------------------------------
\begin{frame}{Mudança de energia ao inserir um dielétrico}

\small

Se as cargas livres que produzem o campo são mantidas fixas, inserir um dielétrico
muda:

\[
\mathbf E_0,\mathbf D_0
\quad \longrightarrow \quad
\mathbf E,\mathbf D.
\]

A mudança de energia é:

\[
\Delta U_E
=
\frac{1}{2}
\int d^3r\,
\left(
\mathbf E\cdot\mathbf D
-
\mathbf E_0\cdot\mathbf D_0
\right).
\]

Depois de rearranjar termos e usar que a carga livre foi mantida fixa, Zangwill
obtém:

\begin{eqbox}
\begin{equation}
\Delta U_E
=
-\frac{1}{2}
\int d^3r\,\mathbf P\cdot\mathbf E_0.
\tag{6.101}
\label{eq:energia-insercao-dieletrico}
\end{equation}
\end{eqbox}

\end{frame}

%-------------------------------------------------
\begin{frame}{De onde vem o fator \(1/2\)?}

\footnotesize

Para um pequeno objeto linear:

\[
\mathbf p=\alpha\epsilon_0\mathbf E_0.
\]

Se o dipolo já estivesse pronto, sua energia de interação seria:

\[
U_{\rm ext}=-\mathbf p\cdot\mathbf E_0.
\]

Mas o dipolo precisa ser criado. Há um custo interno, que para resposta linear é
quadrático em \(p\):

\[
\Delta U_E=-pE_0+\zeta p^2.
\]

Minimizando:

\[
\frac{d}{dp}(-pE_0+\zeta p^2)=0
\quad\Rightarrow\quad
p=\frac{E_0}{2\zeta}.
\]

Como \(p=\alpha\epsilon_0E_0\), o mínimo dá:

\begin{eqbox}
\[
\Delta U_E=-\frac{1}{2}\mathbf p\cdot\mathbf E_0.
\]
\end{eqbox}

\end{frame}

%-------------------------------------------------
\begin{frame}{Transformada de Legendre: potenciais fixos}

\small

Quando condutores estão ligados a baterias, os potenciais \(\varphi_k\) são
mantidos fixos. Como em condutores, usamos uma energia transformada:

\begin{eqbox}
\begin{equation}
\widehat U_E
=
U_E
-
\sum_k Q_k\varphi_k.
\tag{6.109}
\label{eq:Uhat-dieletrico}
\end{equation}
\end{eqbox}

\begin{mybox}{Ideia}
\(U_E\) é natural em cargas livres. \(\widehat U_E\) é natural em potenciais fixos.
\end{mybox}

Para dielétrico simples, essa função é usada para calcular forças quando a
geometria muda mantendo \(\varphi_k\) constantes.

\end{frame}

%=================================================
\sectioncover{6.8 Forças em matéria dielétrica}

%-------------------------------------------------
\begin{frame}{Força sobre um corpo dielétrico isolado}

\small

Para um dielétrico polarizado \(\mathbf P(\mathbf r)\) em campo externo
\(\mathbf E_0(\mathbf r)\), a densidade de carga de polarização sente força:

\[
d\mathbf F=\rho_P\mathbf E_0\,d^3r+\sigma_P\mathbf E_0\,dS.
\]

Assim:

\begin{eqbox}
\begin{equation}
\mathbf F
=
\int_V d^3r\,\rho_P\mathbf E_0
+
\int_S dS\,\sigma_P\mathbf E_0.
\tag{6.115}
\label{eq:forca-coulomb-dieletrico}
\end{equation}
\end{eqbox}

\begin{alertblock}{Importante}
Essa fórmula vale para um corpo isolado em campo externo, não para um subvolume
arbitrário dentro de um dielétrico contínuo.
\end{alertblock}

\end{frame}

%-------------------------------------------------
\begin{frame}{Forma de dipolos}

\small

Como matéria polarizada é equivalente a dipolos com densidade \(\mathbf P\), a
força também pode ser escrita como:

\begin{eqbox}
\begin{equation}
\mathbf F
=
\int_V d^3r\,
(\mathbf P\cdot\nabla)\mathbf E_0.
\tag{6.116}
\label{eq:forca-dipolos-dieletrico}
\end{equation}
\end{eqbox}

\begin{mybox}{Para campo quase uniforme}
Se o corpo é pequeno:
\[
\mathbf F\simeq(\mathbf p\cdot\nabla)\mathbf E_0.
\]
\end{mybox}

\begin{alertblock}{Dieletroforese}
Um dielétrico neutro pode ser atraído para regiões de campo mais intenso.
\end{alertblock}

\end{frame}

%-------------------------------------------------
\begin{frame}{Por que subvolumes são sutis?}

\small

Se cortamos mentalmente um subvolume de um dielétrico contínuo, aparecem forças
microscópicas de curto alcance nas fronteiras artificiais.

\begin{mybox}{Problema}
A fórmula de Coulomb com \(\rho_P\) e \(\sigma_P\) não conta corretamente todas as
forças internas quando a fronteira não é uma superfície física real.
\end{mybox}

Zangwill mostra que isso é crucial para interfaces dielétricas: uma fórmula
ingênua pode errar até por um fator 2 no limite condutor.

\begin{alertblock}{Estratégia correta}
Para dielétricos simples, use variação de energia ou tensor de tensões.
\end{alertblock}

\end{frame}

%-------------------------------------------------
\begin{frame}{Interface dielétrica: tentativa ingênua}

\small

Considere um capacitor de carga livre fixa \(\sigma_f\), com uma interface plana
entre \(\kappa_1\) e \(\kappa_2\). Se usamos apenas a força sobre a carga de
polarização da interface, obtemos:

\begin{eqbox}
\begin{equation}
\mathbf f_{\rm ing}
=
\frac{\sigma_f^2}{\epsilon_0}
\left(
\frac{1}{\kappa_2}
-
\frac{1}{\kappa_1}
\right)\hat{\mathbf z}.
\tag{6.121}
\end{equation}
\end{eqbox}

\begin{alertblock}{Problema}
No limite \(\kappa_1\to\infty\), o meio 1 deveria se comportar como condutor.
Mas o resultado correto para a pressão entre placas teria um fator \(1/2\).
\end{alertblock}

\begin{mybox}{Diagnóstico}
A força sobre um subvolume dielétrico não é apenas a força Coulombiana sobre
\(\rho_P\) e \(\sigma_P\). Faltam forças microscópicas de curto alcance na fronteira
do subvolume.
\end{mybox}

\end{frame}

%-------------------------------------------------
\begin{frame}{Fórmula de Helmholtz}

\small

Para um dielétrico simples, a força total pode ser obtida variando a energia:

\[
U_E=\frac12\int d^3r\,\frac{|\mathbf D|^2}{\epsilon}.
\]

O resultado local é a fórmula de Helmholtz:

\begin{eqbox}
\begin{equation}
\mathbf F
=
\int d^3r
\left[
\rho_f\mathbf E
-
\frac{1}{2}E^2\nabla\epsilon
\right].
\tag{6.130}
\label{eq:helmholtz-force}
\end{equation}
\end{eqbox}

\begin{itemize}
\item \(\rho_f\mathbf E\): força sobre carga livre.
\item \(-\frac12 E^2\nabla\epsilon\): força onde a permissividade varia.
\end{itemize}

\end{frame}

%-------------------------------------------------
\begin{frame}{Exemplo 6.6: força correta na interface}

\small

Pela fórmula de Helmholtz, perto de uma interface sem carga livre:

\[
\mathbf f
=
-\frac12\int dz\,E^2(z)\frac{d\epsilon}{dz}\,\hat{\mathbf z}.
\]

Como \(E_z\) salta na interface, essa forma é delicada. Mas \(D_z\) é contínuo:

\[
D_z=\sigma_f.
\]

Use \(E=D/\epsilon\):

\[
E^2\,d\epsilon
=
\frac{D^2}{\epsilon^2}d\epsilon
=
-D^2\,d\left(\frac{1}{\epsilon}\right).
\]

Então:

\begin{eqbox}
\begin{equation}
\mathbf f
=
\frac{\sigma_f^2}{2\epsilon_0}
\left(
\frac{1}{\kappa_2}
-
\frac{1}{\kappa_1}
\right)\hat{\mathbf z}.
\end{equation}
\end{eqbox}

\end{frame}

%-------------------------------------------------
\begin{frame}{Cheque: limite condutor}

\footnotesize

Se \(\kappa_1\to\infty\), a interface se comporta como uma placa condutora em
contato com o meio 2:

\[
\mathbf f
\to
\frac{\sigma_f^2}{2\kappa_2\epsilon_0}\hat{\mathbf z}.
\]

\begin{mybox}{Comparação com capacitor}
Para um capacitor preenchido pelo meio 2:
\[
C=\frac{\kappa_2\epsilon_0 A}{d}.
\]
Com carga fixa \(Q=\sigma_f A\), a energia é
\[
U_E=\frac{Q^2}{2C}.
\]
\end{mybox}

A força por área é:

\[
\frac{F}{A}
=
\frac{\sigma_f^2}{2\kappa_2\epsilon_0},
\]

em acordo com a fórmula de Helmholtz.

\end{frame}

%-------------------------------------------------
\begin{frame}{Tensor de tensões em dielétrico simples}

\small

A força também pode ser escrita como uma integral de superfície:

\begin{eqbox}
\begin{equation}
\mathbf F
=
\oint_S dS\,
\left[
(\mathbf D\cdot\hat{\mathbf n})\mathbf E
-
\frac12
(\mathbf D\cdot\mathbf E)\hat{\mathbf n}
\right].
\tag{6.132}
\label{eq:forca-stress-superficie}
\end{equation}
\end{eqbox}

Isso motiva o tensor:

\begin{eqbox}
\begin{equation}
T_{ij}
=
D_iE_j
-
\frac12\delta_{ij}\,\mathbf D\cdot\mathbf E.
\tag{6.133}
\label{eq:tensor-stress-dieletrico}
\end{equation}
\end{eqbox}

\[
F_i=\oint_S dS\,T_{ij}\hat n_j.
\tag{6.134}
\]

\end{frame}

%-------------------------------------------------
\begin{frame}{Exemplo: slab dielétrico em capacitor}

\small

Um slab dielétrico parcialmente inserido em um capacitor é puxado para dentro.

\begin{center}
\begin{tikzpicture}[scale=0.9]
\draw[thick] (-3,1) -- (3,1);
\draw[thick] (-3,-1) -- (3,-1);
\node at (3.35,1) {\(+\)};
\node at (3.35,-1) {\(-\)};
\fill[mygray,draw=myblue,thick] (-3,-0.85) rectangle (-0.4,0.85);
\node at (-1.7,0) {\(\epsilon\)};
\draw[->,very thick,myred] (-0.35,0) -- (0.8,0) node[midway,above] {\(\mathbf F\)};
\foreach \x in {-2.5,-1.7,-0.9,0.1,0.9,1.7,2.5}{
  \draw[->,thin] (\x,0.75) -- (\x,-0.75);
}
\end{tikzpicture}
\end{center}

\begin{mybox}{Energia}
O sistema reduz sua energia ao preencher mais da região de campo forte com
material de maior \(\epsilon\).
\end{mybox}

\end{frame}

%-------------------------------------------------
\begin{frame}{Slab no capacitor: cálculo da força}

\small

Para potencial fixo, a energia relevante é a transformada de Legendre
\(\widehat U_E\). Se o slab avança uma distância \(s\), o volume adicional
preenchido é \(Lds\).

\[
\Delta\widehat U_E
=
-\frac12\epsilon_0(\kappa-1)
\left(\frac{V}{d}\right)^2(Lds).
\]

\begin{mybox}{Sinal}
O sinal negativo diz que a energia diminui quando mais dielétrico entra na
região entre as placas.
\end{mybox}

\begin{eqbox}
\[
\mathbf F
=
-\frac{\partial\Delta\widehat U_E}{\partial s}\hat{\mathbf x}
=
\frac{\epsilon_0(\kappa-1)LV^2}{2d}\hat{\mathbf x}.
\]
\end{eqbox}

\end{frame}

%-------------------------------------------------
\begin{frame}{Slab no capacitor: interpretação}

\small

\begin{alertblock}{Mensagem prática}
Dielétricos são atraídos para regiões onde o campo elétrico é mais intenso.
\end{alertblock}

\begin{mybox}{O que a fórmula está dizendo}
\[
F\propto(\kappa-1),
\qquad
F\propto V^2,
\qquad
F\propto \frac{L}{d}.
\]
O efeito cresce com a polarizabilidade do material, com o campo aplicado e com a
área de borda que entra no capacitor.
\end{mybox}

\end{frame}

%-------------------------------------------------
\begin{frame}{Resumo do capítulo 6}

\footnotesize

\begin{mybox}{Fórmulas essenciais}
\begin{itemize}
\item \(\rho_P=-\nabla\cdot\mathbf P\), \(\sigma_P=\mathbf P\cdot\hat{\mathbf n}\).
\item \(\displaystyle \int_V\mathbf P\,d^3r=\mathbf p\).
\item \(\mathbf D=\epsilon_0\mathbf E+\mathbf P\), \(\nabla\cdot\mathbf D=\rho_f\).
\item Dielétrico simples: \(\mathbf D=\epsilon\mathbf E=\kappa\epsilon_0\mathbf E\).
\item Energia: \(\displaystyle U_E=\int d^3r\int_0^{\mathbf D}\mathbf E\cdot\delta\mathbf D\).
\item Para dielétrico simples: \(\displaystyle U_E=\frac12\int d^3r\,\mathbf E\cdot\mathbf D\).
\end{itemize}
\end{mybox}

\begin{alertblock}{Ideia unificadora}
A resposta dielétrica é uma competição entre energia eletrostática, polarização
induzida e forças microscópicas que mantêm a matéria coesa.
\end{alertblock}

\end{frame}

\end{document}
